\section{Síntesis y ampliación}

\subsection{Repaso del contenido del curso}

El curso KAIST CS492D se centró esta sesión en la teoría y aplicaciones de los modelos de difusión y Flow, cubriendo dos grandes áreas:

\subsubsection{Fundamentos teóricos}

\begin{enumerate}
    \item \textbf{DDPM} (Denoising Diffusion Probabilistic Models): marco base de los modelos de difusión, definiendo el proceso de adición de ruido hacia adelante y el proceso de eliminación de ruido hacia atrás.
    \item \textbf{DDIM} (Denoising Diffusion Implicit Models): método de muestreo determinista que acelera el proceso de inferencia.
    \item \textbf{Difusión en tiempo continuo}: extensión del proceso de difusión discreto al marco de EDO/SDE en tiempo continuo.
    \item \textbf{Conversión entre SDE y ODE}: establece la conexión entre muestreo estocástico y determinista.
    \item \textbf{Flow Matching}: método más flexible de modelado generativo que conecta los modelos de difusión con la teoría de transporte óptimo.
\end{enumerate}

\subsubsection{Expansión de aplicaciones}

\begin{enumerate}
    \item \textbf{Guía en tiempo de inferencia} (Inference-Time Guidance): control del proceso generativo mediante señales de guía sin necesidad de volver a entrenar el modelo.
    \item \textbf{Destilación de puntuaciones} (Score Distillation): transferencia del conocimiento de modelos de difusión preentrenados a otros dominios de datos, logrando generación cruzada de modalidades.
    \item \textbf{Generación 3D}: creación de contenido 3D aprovechando conocimientos previos de modelos de difusión de imágenes 2D.
\end{enumerate}

\begin{importantbox}{Filosofía del diseño del curso}
El diseño del curso para este semestre enfatiza un enfoque progresivo desde los fundamentos hasta las aplicaciones: primero se comprende la esencia matemática de los modelos (teoría de probabilidades, ecuaciones diferenciales estocásticas, transporte óptimo), y luego se explora cómo aplicar estas herramientas teóricas a problemas prácticos (edición de imágenes, generación 3D, transferencia cruzada de modalidades).
\end{importantbox}

\subsection{Futuras direcciones: difusión discreta y modelos multimodales}

Al final del curso, la profesora Minhyuk Sung planteó una pregunta reflexiva:

\begin{importantbox}{Autoregresivo vs. Difusión: el límite se difumina}
La visión tradicional sostenía que:
\begin{itemize}
    \item Los \textbf{modelos autoregresivos} son ideales para datos de secuencia \textbf{discreta} como el texto.
    \item Los \textbf{modelos de difusión/Flow} son ideales para datos de alta dimensión \textbf{continuos} como imágenes y videos.
\end{itemize}

Sin embargo, este límite está rompiéndose:
\begin{itemize}
    \item Los modelos autoregresivos también generan imágenes de alta calidad mediante tokenización de imágenes (por ejemplo, VQ-VAE).
    \item Los modelos de difusión discreta extienden el proceso de difusión al espacio de datos discretos.
    \item Modelos multimodales como GPT-4o y Gemini pueden generar texto e imágenes simultáneamente.
\end{itemize}
\end{importantbox}

\subsubsection{Modelos de difusión discretos}

Los modelos de difusión no se limitan a datos continuos; también pueden aplicarse a datos discretos. El proceso hacia adelante de la difusión discreta no implica agregar ruido gaussiano, sino degradar gradualmente los datos mediante una \textbf{cadena de Markov discreta} (por ejemplo, reemplazo aleatorio de tokens). Esto abre la puerta para aplicar modelos de difusión al texto, código y secuencias moleculares.

\subsubsection{Generación multimodal unificada}

La investigación de vanguardia actual explora unificar el texto y las imágenes en un mismo marco generativo:

\begin{knowledgebox}{Dos paradigmas para arquitecturas multimodales unificadas}
\begin{enumerate}
    \item \textbf{Tokenizar todo}: discretización de imágenes en tokens (mediante VQ-VAE, etc.) y procesamiento conjunto con tokens de texto mediante modelos autoregresivos. Ejemplo representativo: la serie GPT-4o.
    \item \textbf{Arquitectura híbrida}: decodificación autoregresiva para la parte de texto y decodificación por difusión/Flow para la parte de imágenes, conectadas a través de un espacio de representación compartido. Ejemplo representativo: Transfusion.
\end{enumerate}

Cuál paradigma prevalecerá sigue siendo una pregunta abierta.
\end{knowledgebox}

\subsection{Sugerencias para investigadores}

La profesora Minhyuk Sung ofreció varias recomendaciones al finalizar el curso:

\begin{enumerate}
    \item \textbf{Practicar}: existen numerosas implementaciones de código abierto para técnicas como SDS; se recomienda ejecutar personalmente el código, observar los resultados y comprender los patrones de falla.
    \item \textbf{Enfocarse no solo en la generación, sino también en la edición}: la edición 3D (modificación de contenido 3D existente) puede tener un valor práctico mayor que la generación ex nihilo.
    \item \textbf>Pensar transversalmente>: el marco de SDS no se limita al contenido visual; existen aplicaciones potenciales en audio, moléculas y movimiento.
    \item \textbf{Seguir tendencias de vanguardia}: las áreas de fusión entre autoregresivo y difusión, así como modelos multimodales unificados, merecen atención continua.
\end{enumerate}

\subsection{Convertir el contenido de esta clase en un plan de investigación}

Si se considera esta clase como un punto de partida para la investigación en lugar de una simple revisión, lo más importante es descomponer aún más ``¿qué puede hacer SDS'' en rutas experimentales ejecutables. Muchos proyectos de texto a 3D fracasan no por falta de conocimiento previo de modelos grandes, sino porque no se diseñan claramente las tres cosas: \textbf{representación, optimización y evaluación}.

\begin{knowledgebox}{Diseño experimental mínimo para proyectos SDS}
\begin{table}[H]
\centering
\begin{tabular}{p{0.16\textwidth} p{0.28\textwidth} p{0.34\textwidth}}
\toprule
Módulo & Selección mínima viable & Observaciones que deben registrarse prioritariamente \\
\midrule
Representación 3D & NeRF o 3D Gaussian Splatting & Coherencia de perspectiva, suavidad geométrica, estabilidad del entrenamiento \\
Priorio 2D & Modelos de difusión condicionados por texto tipo Stable Diffusion / Imagen & Sensibilidad al prompt, saturación de color, recuperación de detalles \\
Estrategia de optimización & Baseline SDS + regularización + muestreo de cámara & Frecuencia Janus, deriva de parches, velocidad de convergencia \\
Método de evaluación & Renderizado multivista + puntuación CLIP/usuario + imprimibilidad geométrica & ¿El resultado solo ``parece'' o tiene verdadera consistencia 3D? \\
\bottomrule
\end{tabular}
\caption{Construcción del primer prototipo de investigación de texto a 3D basado en el contenido del curso}
\end{table}
\end{knowledgebox}

\begin{warningbox}{Juicio erróneo común: alta calidad 2D no implica alta calidad 3D}
Una imagen renderizada individualmente hermosa no demuestra que el modelo haya aprendido realmente geometría estable. El error de juicio más común en el sistema SDS ocurre cuando los investigadores se convence por imágenes de alta fidelidad en ciertos puntos de vista, ignorando problemas como colapso estructural al girar a la parte trasera, texturas repetidas y errores topológicos del objeto. Por lo tanto, los registros experimentales deben incluir una crucero continuo de perspectivas, no solo muestras óptimas.
\end{warningbox}

\subsection{Patrones de falla y puntos de control para ajuste de parámetros}

Desde la perspectiva de la práctica investigativa, la parte más difícil de un proyecto SDS no es ``ejecutar el loss'', sino identificar en qué nivel surge exactamente el fallo actual: si las condiciones de texto son insuficientemente claras, si el priorio de difusión 2D es demasiado fuerte, si el muestreo de cámara es irrazonable o si la capacidad del propio sistema de representación 3D es limitada. Mezclar estos problemas lleva a que los experimentos queden reducidos a probar prompts ciegamente.

\begin{table}[H]
\centering
\begin{tabular}{p{0.22\textwidth} p{0.3\textwidth} p{0.28\textwidth}}
\toprule
Fenómeno & Causa raíz más probable & Acción de reparación a intentar prioritariamente \\
\midrule
Adecuado por delante, colapsado por detrás & Restricciones multivista insuficientes, distribución de perspectiva de muestreo demasiado estrecha & Ampliar la cobertura de elevación/azimut, introducir regularización dependiente de la vista \\
Color sobresaturado, texturas pegajosas & Gradiente SDS demasiado fuerte, arrastre excesivo del priorio 2D & Reducir la intensidad de guía, aumentar entropía/regularización TV \\
Geometría hinchada, bordes difusos & Capacidad de representación insuficiente o regularización de densidad inestable & Cambiar sistema de representación 3D, agregar priorio geométrico, restringir explícitamente opacidad \\
Resultados altamente dependientes de detalles del prompt & Las condiciones de texto asumen demasiada información estructural & Reescribir plantillas de prompts y acompañar con perspectivas de referencia o supervisión débil \\
\bottomrule
\end{tabular}
\caption{Patrones de falla comunes y acciones de reparación en proyectos SDS de texto a 3D}
\end{table}

\begin{importantbox}{El verdadero marco metodológico dejado por el cierre del curso}
La lección más importante al final de este curso no es ``aprender otro algoritmo nuevo'', sino comprender que la investigación sobre modelos de difusión está pasando de la mera generación de imágenes a \textbf{transferir priores fuertes a dominios con datos escasos}. SDS, VSD, difusión discreta y arquitecturas multimodales unificadas pertenecen a este problema matriz más amplio: cómo permitir que un generador entrenado en conjuntos de datos masivos proporcione señales de optimización para otro espacio de representación valioso pero con pocos datos.
\end{importantbox}

\subsection{Resumen de la sección}

El curso cubrió sistemáticamente el desarrollo teórico desde DDPM hasta Flow Matching, así como la expansión de aplicaciones desde la guía en tiempo de inferencia hasta la destilación de puntuaciones. Las tendencias de vanguardia actuales —difusión discreta, generación de imágenes autoregresiva y modelos multimodales unificados— están rompiendo los límites tradicionales de selección de modelos, presagiando marcos de modelado generativo más flexibles y potentes.

\newpage
